\chapter[纵然那是纯粹道德之举]{纵然那是纯粹道德之举}

\chaptersubtitle{《这种事绝对很奇怪啊》}

\chapternote{%
    ``\textbf{Even a purely moral act} that has no hope of any immediate and visible political effect can gradually and indirectly, over time, gain in political significance.''
    \par
    {\raggedleft\noindent ---Václav Havel, \emph{Letter to Alexander Dubček}\par}
    \par
    “纵然那是纯粹道德之举眼下并不产生任何直接或明显的政治效果，在长期的历史进程中，它仍可能酝酿出深远的政治意义。”
    \par
    {\raggedleft\noindent ——瓦茨拉夫·哈维尔《致亚历山大·杜布切克的信》\par}
    \par
    译者按：第六集揭开灵魂宝石真相的那一刻，私人愿望的代价，第一次露出了它被安排的面目。纵是看似私人的道德之举，也能在漫长的暗流中显出制度的轮廓。
}

有位大学教授曾给过我这样一条颇为古怪的建议：
如果你在写作时卡壳了，不妨找到作品主题的确切中点，然后围绕你在那里发现的任何内容展开论述。
虽然我在讨论《魔圆》的主题时并不担心灵感枯竭，但既然本集末尾恰好是整部剧的中点，此刻来讨论剧中涉及的“同意权”问题，再合适不过了。

本集最后一幕中，角色们因丘比揭露的真相而惊恐万分：
魔法少女的身体并非活体，而仅仅是躯壳；
若灵魂宝石完好无损，身体便能无限修复。
唯有破坏灵魂宝石，才能真正伤害魔法少女本身——
而小圆将沙耶香的灵魂宝石扔下桥，实际上就导致了沙耶香的暂时性死亡（多亏焰及时将它找回）。

诚然，不妨说（尽管这论点站不住脚，但总归能提一提）：
女孩们的激烈反应纯属小题大做。
坦白讲，丘比所描述的交易听起来相当诱人：
身体的生理体验与活着时几乎无异，大多数魔法少女甚至从未察觉变化，而且这具身体不仅十分健康，更为耐用，还能治愈任何伤病。
此外，从杏子的情况来看，这副身体似乎还能毫无负担地享用垃圾食品。
换作是我，那肯定会立刻答应。

事实上，灵魂宝石的设定显然借鉴了俄罗斯民间传说角色 \emph{Koshchyei Byessmyertnuy}（不死者科西切）。
他将自己的灵魂藏在手指里，切断手指后将其藏在一枚蛋中，蛋藏在一只鸭子里，鸭子藏在一只兔子中，兔子又被装进一个铁箱里，最后铁箱被埋在遥远岛屿上的一棵绿橡树下。
科西切是一位诱拐少女的恶棍，只有当英雄找到那个蛋时，才能伤害到他\cite{frazer1982external-soul}。
科西切这样做的优势显而易见，这很可能也是流行文化中他的“后继者”采用类似保护手段的原因（其中最广为人知的例子当属《龙与地下城》\cite{wilkes2003lich}中巫妖的命匣和《哈利·波特》\cite{rowling2005half-blood-prince}中伏地魔的魂器）。

然而，科西切与沙耶香之间的根本区别，在于是否获得了知情且明确的同意。
传说表明，科西切知晓自己在做什么并主动选择这么做。
而沙耶香却全然不知自己的生命已系于灵魂宝石，更不知道自己的身体已在违背其意愿的情况下被强行改造。
在后续剧情中，她提到自己不相信这具新身体能够生育——
这或许曾是她对未来生活的美好期许。
无论如何，沙耶香、杏子和小圆在本集中表现出的恐惧充分表明，三人都将此视为一种根本性的侵犯。

丘比辩解说，它不理解人类为何如此在意灵魂的位置。
但这纯属避重就轻的胡扯，既然它心知肚明人类对此十分在意，纠结其原因本身就毫无意义。
丘比在签订契约时故意隐瞒关键信息，随后又在受害者无法接受真相时指责她们。
本质上，它是在用“沙耶香从未说过‘不’”来为自己的行为开脱。

“不就是不”在争取女性权利、尤其是涉及同意与身体自主权的运动中，常被用作口号。
然而，丘比的例子表明，虽然承认“不就是不”肯定比不承认要好，但这仍是一个不完整的标准。
比“不就是不”更重要的是“是就是是”，这正是“肯定性同意”的核心含义\cite{caringella2009affirmative-consent}。
沉默并不能为侵犯行为提供正当理由，因为这可能意味着当事人无法反对，正如沙耶香无法反对她根本不知情的契约条款一样。

如何尊重他人的选择与自主权，在本集另一幕中得到了耐人寻味的呼应：
小圆（自然使用了最含糊的措辞）与母亲谈起沙耶香的近况。
这里有两点至关重要：
第一，母亲指出，做正确的事并不总能带来幸福或好结果。
这实质上明确拒斥了结果主义伦理立场（一种伦理学观点，认为行为的道德性取决于其后果\cite{sinnott-armstrong-nodate-consequentialism}），与整部剧所持的立场大体一致（正因如此，丘比始终被刻画为一个纯粹的结果主义者）。

在与母亲的对话场景中，动画明确拒绝结果主义，其深意延伸至本集结尾的桥头场景——
后者同样以隐晦方式否定了该立场。
丘比的立场是结果主义的：
灵魂抽取对魔法少女有益，因为她们得以对抗魔女并活下来，但得知此事往往令她们痛苦不堪，因此最佳做法就是抽取灵魂但不告知真相。
该场景的音乐与构图（尤其体现在：丘比虽然体型小巧，但画面却让它的阴影笼罩着女孩们）清晰地表明了，剧作不仅拒绝丘比的说辞，更与女孩们的恐惧共情，从而彻底摒弃结果主义观点。

这场母女对谈的第二个深意在于，将沙耶香刻画为“秉正心而为，招悖逆之果”的角色——
这一评价绝非她所独有，也适用于焰。
她反复尝试拯救小圆，结果却让小圆在每一次时间轮回中承受更多痛苦，最后成为了越发强大的魔女。
母亲的建议也促使小圆为朋友而行差踏错——
这不仅指将沙耶香的灵魂宝石扔下桥，也指她在最终话选择成为魔法少女（这正是焰一直竭力阻止的事）。
鉴于背景音乐的选择，这种双重含义并非偶然：
本集母女对谈的配乐 \emph{Clementia}（怜悯），在最终话小圆许愿后的 \emph{Sagitta Luminis}（光之箭矢）中以变奏形式再现。

此外，这种“为拯救他人而犯错”的理念在《叛逆》中也得到了重申，小圆与焰都为对方采取了看似极不明智的行动——
不过更多细节，留待到那时再展开。
